% Session 2: Research Pipelines — From Data to Results
% Beamer Presentation — Claude Code for Academic Researchers
\documentclass[aspectratio=169, 11pt]{beamer}

% ── Theme and Colors ──────────────────────────────────────────────
\usetheme{Madrid}
\usecolortheme{default}

% Custom colors
\definecolor{anthropic}{RGB}{50, 50, 80}
\definecolor{accent}{RGB}{180, 90, 40}
\definecolor{lightgray}{RGB}{245, 245, 245}

\setbeamercolor{structure}{fg=anthropic}
\setbeamercolor{frametitle}{bg=anthropic, fg=white}
\setbeamercolor{title}{bg=anthropic, fg=white}
\setbeamercolor{block title}{bg=anthropic, fg=white}
\setbeamercolor{block body}{bg=lightgray}
\setbeamercolor{alerted text}{fg=accent}

\setbeamerfont{title}{size=\Large, series=\bfseries}
\setbeamerfont{frametitle}{size=\large, series=\bfseries}

% ── Packages ──────────────────────────────────────────────────────
\usepackage{graphicx}
\usepackage{booktabs}
\usepackage{listings}
\usepackage{fontenc}
\usepackage{hyperref}
\usepackage{tikz}
\usetikzlibrary{positioning}

\lstset{
  basicstyle=\ttfamily\small,
  backgroundcolor=\color{lightgray},
  frame=single,
  framerule=0pt,
  breaklines=true,
  columns=fullflexible,
  literate={_}{\_}1,
}

% ── Metadata ──────────────────────────────────────────────────────
\title{Session 2: Research Pipelines}
\subtitle{From Data to Results}
\author{Brady Allardice}
\date{}
\institute{UPF \& IBEI}

% ══════════════════════════════════════════════════════════════════
\begin{document}

% ── Title ─────────────────────────────────────────────────────────
\begin{frame}[plain]
  \titlepage
\end{frame}

% ══════════════════════════════════════════════════════════════════
%  PART 1: THE CORE WORKFLOW
% ══════════════════════════════════════════════════════════════════

\begin{frame}{Where We Are}
  \textbf{Session 1:} What Claude Code can do. How it works. CLAUDE.md and plan mode.

  \vspace{1em}
  \textbf{Session 2:} How to use it on a real research task --- and stay in control.

  \vspace{1em}
  Today you will:
  \begin{enumerate}
    \item Merge two datasets with explicit instructions
    \item Estimate a model and verify the output
    \item Generate robustness checks
    \item Produce a \LaTeX{} results table
  \end{enumerate}

  \vspace{0.5em}
  \alert{The methods are straightforward. The skill is the workflow.}
\end{frame}

% ══════════════════════════════════════════════════════════════════
%  STRATEGIES FOR WORKING WITHIN THE CONTEXT WINDOW
% ══════════════════════════════════════════════════════════════════

\begin{frame}{Strategies for Working Within the Context Window}
  Recall from Session 1: Claude Code can only ``see'' a limited amount of text at once.\\
  Errors compound when it loses track of earlier decisions.

  \vspace{0.5em}
  \textbf{Five strategies to work within this constraint:}
  \begin{enumerate}
    \item Keep tasks focused
    \item Maintain project documentation
    \item Summarize work between conversations
    \item Validate iteratively
    \item Write token-aware prompts
  \end{enumerate}
\end{frame}

\begin{frame}{What Happens When the Context Window Fills Up}
  As a conversation grows, Claude Code eventually runs out of room.\\
  When this happens, it \textbf{auto-compacts}: it summarizes older parts of the conversation to free up space.

  \vspace{0.5em}
  \textbf{What auto-compacting does:}
  \begin{itemize}
    \item Replaces earlier messages with a compressed summary
    \item Keeps recent messages and tool results intact
    \item Happens automatically. You will see a notice in the chat
  \end{itemize}

  \vspace{0.5em}
  \textbf{The risk:} Details from early in the conversation (variable definitions, decisions you discussed, constraints you mentioned) may be lost or simplified in the summary.

  \vspace{0.5em}
  \textbf{You can also compact manually:}\\
  Type \texttt{/compact} to force a summary and free up context space.

  \vspace{0.5em}
  \alert{This is why CLAUDE.md matters:} it survives compacting because it is reloaded from disk, not from conversation history.
\end{frame}

\begin{frame}{Strategy 1: Keep Tasks Focused}
  \centering
  {\large \textbf{One topic, one conversation.}}

  \vspace{1em}
  \begin{columns}[T]
    \begin{column}{0.48\textwidth}
      \textbf{\color{red!70!black} Avoid:}
      \begin{itemize}
        \item ``Clean the data, build the panel, run the regressions, and make tables''
        \item Long conversations that drift across topics
        \item Asking Claude Code to hold your entire project in its head
      \end{itemize}
    \end{column}
    \begin{column}{0.48\textwidth}
      \textbf{\color{green!50!black} Instead:}
      \begin{itemize}
        \item ``Clean and reshape the raw census file''
        \item Start a new conversation for each distinct task
        \item Let each conversation do one thing well
      \end{itemize}
    \end{column}
  \end{columns}

  \vspace{0.8em}
  \begin{alertblock}{Why}
    A focused task fits comfortably in the context window. A sprawling task forces Claude Code to forget earlier steps, and that is where silent errors enter.
  \end{alertblock}
\end{frame}

\begin{frame}{Strategy 2: Maintain Project Documentation}
  Claude Code reads certain files \textbf{automatically} every time it starts.

  \vspace{0.5em}
  \textbf{Keep these up to date:}
  \begin{itemize}
    \item \texttt{CLAUDE.md}: conventions, constraints, variable definitions
    \item \texttt{README.md}: project structure, what each directory contains
    \item Directory-level \texttt{CLAUDE.md} files for subfolders with specific rules
  \end{itemize}

  \vspace{0.5em}
  \textbf{Tell Claude Code what to \textit{ignore}:}
  \begin{itemize}
    \item \texttt{.claudeignore} works like \texttt{.gitignore}
    \item Exclude large data files, build artifacts, logs
    \item Prevents Claude Code from wasting context on irrelevant files
  \end{itemize}

  \vspace{0.5em}
  \begin{block}{The Payoff}
    Good documentation means Claude Code starts every conversation already knowing your conventions, without you spending tokens re-explaining them.
  \end{block}
\end{frame}

\begin{frame}[fragile]{Strategy 3: Summarize Work Between Conversations}
  When you finish a conversation, the context is gone. The next conversation starts fresh.

  \vspace{0.8em}
  \textbf{Before ending a conversation:}
  \begin{enumerate}
    \item Ask Claude Code to summarize what it did and what remains
    \item Save that summary (in a file, a commit message, or your notes)
    \item Paste the summary into the \textit{next} conversation as context
  \end{enumerate}

  \vspace{0.8em}
  \textbf{Example prompt:}
  \begin{lstlisting}
Summarize what we did in this session: what files
were changed, what decisions were made, and what
still needs to be done.
  \end{lstlisting}

  \vspace{0.5em}
  \alert{Think of it like a handoff memo between research assistants.}\\
  The next ``assistant'' knows nothing unless you brief it.
\end{frame}

\begin{frame}{Strategy 4: Validate Iteratively}
  \begin{alertblock}{The Compounding Error Problem}
    If step 1 is wrong and you do not catch it, steps 2--5 build on that mistake. By the time you check the final output, the error is buried.
  \end{alertblock}

  \vspace{0.5em}
  \textbf{Check each step before moving to the next:}
  \begin{enumerate}
    \item Clean the data $\rightarrow$ \textit{verify row counts and distributions}
    \item Construct variables $\rightarrow$ \textit{spot-check values against known cases}
    \item Run the merge $\rightarrow$ \textit{confirm pre- and post-merge row counts}
  \end{enumerate}

  \vspace{0.5em}
  Slower per step, but \textbf{much faster overall} because you never have to unwind a chain of errors you didn't catch.
\end{frame}

\begin{frame}{Strategy 5: Token-Aware Prompting}
  Every word in your prompt costs tokens. Make them count.

  \vspace{0.5em}
  \textbf{A good prompt tells Claude Code three things:}
  \begin{enumerate}
    \item \textbf{What files matter:} ``Read \texttt{data/raw/census.csv} and \texttt{scripts/clean.py}''
    \item \textbf{What the goal is:} ``Create a county-year panel with one row per county per decade''
    \item \textbf{What constraints apply:} ``Keep all 3,143 counties; do not drop rows with missing population''
  \end{enumerate}

  \vspace{0.5em}
  \begin{columns}[T]
    \begin{column}{0.48\textwidth}
      \textbf{\color{red!70!black} Vague:}\\
      {\small ``Clean up the data and make it ready for analysis''}
    \end{column}
    \begin{column}{0.48\textwidth}
      \textbf{\color{green!50!black} Specific:}\\
      {\small ``Read \texttt{census.csv}, drop rows where \texttt{year < 1950}, and save as \texttt{panel\_clean.csv}''}
    \end{column}
  \end{columns}

  \vspace{0.8em}
  \begin{block}{The Rule}
    Be specific about \textit{inputs}, \textit{outputs}, and \textit{constraints}. Be brief about everything else. Context is finite; spend it on what matters.
  \end{block}
\end{frame}

% ══════════════════════════════════════════════════════════════════
%  THE CORE WORKFLOW
% ══════════════════════════════════════════════════════════════════

% ── Core Workflow ──────────────────────────────────────────────────
\begin{frame}{The Core Workflow}
  \begin{tabular}{clll}
    & \textbf{Step} & \textbf{Who} & \textbf{If something is wrong} \\[0.2em]
    \toprule
    1. & Describe & You & Rethink before sending \\[0.3em]
    2. & Review plan & You \textcolor{accent}{\textbf{*}} & Reject, re-instruct \\[0.3em]
    3. & Generate \& validate & Claude Code + You & Read the code, validate, re-instruct \\[0.3em]
    4. & Save or re-instruct & You & Keep the result (git in Session 3) \\
    \bottomrule
  \end{tabular}

  \vspace{0.8em}
  \textcolor{accent}{\textbf{*}} = checkpoint where you can catch a problem before anything gets saved.

  \vspace{0.3em}
  Validation is not a separate step --- it happens \textbf{continuously} throughout steps 1--3.

  \vspace{0.3em}
  \begin{block}{The pattern}
    When something goes wrong, go back to \textbf{whichever step actually fixes the problem}: a bad plan means rewriting the instruction (step 1), not just re-running the same one. Diagnose first, then go back to the right step.
  \end{block}
\end{frame}

\begin{frame}{Step 1: Describe --- Write It Down First}
  \textbf{Before you type anything into Claude Code, write down your instruction.}\\
  If you cannot state the task precisely on paper, Claude Code cannot execute it precisely in code.

  \vspace{0.5em}
  \textbf{A good instruction specifies three things:}
  \begin{enumerate}
    \item \textbf{Inputs:} ``Read \texttt{data/survey.csv} and \texttt{data/demographics.csv}''
    \item \textbf{Goal:} ``Left-join demographics onto the survey by respondent ID''
    \item \textbf{Constraints:} ``Output must have exactly 2,044 rows. Do not drop respondents.''
  \end{enumerate}

  \vspace{0.5em}
  \begin{columns}[T]
    \begin{column}{0.48\textwidth}
      \textbf{\color{red!70!black} Vague:}\\
      {\small ``Merge the datasets and run a regression''}
    \end{column}
    \begin{column}{0.48\textwidth}
      \textbf{\color{green!50!black} Specific:}\\
      {\small ``Left-join demographics onto the survey by respondent\_id. Use control as the reference. Run a weighted logistic regression.''}
    \end{column}
  \end{columns}

  \vspace{0.5em}
  \alert{Be specific about inputs, outputs, and constraints. Be brief about everything else.}
\end{frame}

\begin{frame}{Step 2: Review the Plan}
  In plan mode, Claude Code proposes an approach before acting.

  \vspace{0.5em}
  \textbf{What to look for:}
  \begin{itemize}
    \item Does it read the right files?
    \item Does the merge/analysis strategy match what you specified?
    \item Does it handle edge cases (missing data, extra rows, reference categories)?
    \item Does it make any \textit{analytical decisions} you did not ask for?
  \end{itemize}

  \vspace{0.5em}
  \begin{alertblock}{The key question}
    Is there anything in this plan that I did not explicitly ask for?\\
    If yes: reject or modify before it runs.
  \end{alertblock}

  \vspace{0.5em}
  \textbf{You do not always need plan mode.} For simple tasks --- summary statistics, reformatting, boilerplate --- steps 1--3 collapse into one. But when uncertain, \textit{use it}.
\end{frame}

\begin{frame}[fragile]{Step 3: Generate and Validate}
  Claude Code writes and runs the code. Your job is not done.

  \vspace{0.5em}
  \textbf{Read what Claude Code produced:}
  \begin{itemize}
    \item Does the script do what the plan said it would?
    \item Are variable names and file paths correct?
    \item Did it make any choices the plan did not mention?
  \end{itemize}

  \vspace{0.5em}
  \textbf{Common surprises at this stage:}
  \begin{itemize}
    \item Silently dropping rows with missing values
    \item Choosing a default reference category you did not specify
    \item Using unweighted estimation when you asked for weighted
    \item Adding ``helpful'' data cleaning steps you did not request
  \end{itemize}

  \vspace{0.5em}
  \begin{block}{The habit}
    Read the code, not just the output. Output can look correct even when the code is wrong.
  \end{block}
\end{frame}

% ══════════════════════════════════════════════════════════════════
%  VALIDATION STRATEGIES
% ══════════════════════════════════════════════════════════════════

\begin{frame}{Validation Strategy 1: Read the Code}
  Claude Code is not deterministic --- the same prompt can produce different code each time.\\
  But \textbf{code itself is deterministic.} If you read the code and understand what it does, you can be confident in the results, the same way you would be with code you wrote yourself.

  \vspace{0.8em}
  \textbf{This is the most fundamental verification:}
  \begin{itemize}
    \item Read the script Claude Code produced, line by line
    \item Confirm it does what you asked --- not more, not less
    \item If you understand the code and it is correct, the output is correct
  \end{itemize}

  \vspace{0.8em}
  \begin{block}{The principle}
    You are not trusting Claude Code. You are trusting \textit{the code} --- the same way you trust code you wrote yourself. The difference is who typed it, not whether it works.
  \end{block}
\end{frame}

\begin{frame}{Validation Strategy 2: Check Known Properties}
  Before you look at results, ask: \textit{what do I already know must be true?}

  \vspace{0.5em}
  \textbf{Examples:}
  \begin{itemize}
    \item A left join should not change the number of rows in the left table
    \item Adding a control variable should not change the sample size (unless it has missing values)
    \item A share variable must be between 0 and 1
    \item Treatment groups should sum to the total sample
  \end{itemize}

  \vspace{0.5em}
  \textbf{Ask Claude Code to assert these explicitly:}
  \begin{itemize}
    \item ``Assert the merged output has exactly 2,044 rows''
    \item ``Print the number of observations used in each regression''
    \item ``Verify no duplicate respondent IDs''
  \end{itemize}

  \vspace{0.5em}
  If the assertion fails, you catch the error immediately --- not three steps later.
\end{frame}

\begin{frame}{Validation Strategy 3: Write Down Expectations First}
  \begin{block}{The Pre-Registration Analogy}
    Just as you pre-register hypotheses before seeing results, write down what a correct transformation should produce \textit{before} you run it.
  \end{block}

  \vspace{0.5em}
  \textbf{Before a merge or transformation, state:}
  \begin{itemize}
    \item Expected output dimensions (rows $\times$ columns)
    \item Which columns should appear and their types
    \item How missing values should be handled
    \item A specific row you can trace through by hand
  \end{itemize}

  \vspace{0.5em}
  \textbf{Then compare the actual result to your expectations.}\\
  Discrepancies are either a bug or a misunderstanding --- both worth catching early.
\end{frame}

\begin{frame}[fragile]{Validation Strategy 4: Unit Tests Against Known Values}
  Pick a handful of observations and compute the correct answer by hand.\\
  Then verify the code reproduces them.

  \vspace{0.5em}
  \textbf{Example prompt:}
  \begin{lstlisting}
Write a test that verifies:
- Respondent 9 (interviewed 2015-10-07) has
  treatment = "history" and fx_status = "none"
- The control group has exactly 480 respondents
- The weighted proportion supporting intervention
  is approximately 0.486
  \end{lstlisting}

  \vspace{0.5em}
  \textbf{Why this works:}
  \begin{itemize}
    \item Forces you to engage with the data at the observation level
    \item Catches silent recoding, wrong joins, and off-by-one errors
    \item Tests survive across refactors --- rerun them after every change
  \end{itemize}
\end{frame}

\begin{frame}{Validation Strategy 5: Replicate Before You Extend}
  \textbf{Building on someone else's work?}\\
  Have Claude Code reproduce their results first.

  \vspace{0.5em}
  \begin{enumerate}
    \item Give Claude Code the original paper's code or description
    \item Ask it to replicate a key table or figure
    \item Compare output to the published result
    \item Only after replication succeeds, begin your extension
  \end{enumerate}

  \vspace{0.5em}
  \begin{alertblock}{Why}
    If you cannot reproduce the baseline, anything you build on top of it is uninterpretable. Replication is not a detour; it is the foundation.
  \end{alertblock}
\end{frame}

\begin{frame}{Validation Strategy 6: Use a Second Agent as a Critic}
  Use a \textit{second} Claude Code conversation as a reviewer.

  \vspace{0.5em}
  \textbf{How:}
  \begin{enumerate}
    \item Agent 1 writes the code
    \item Open a new conversation (Agent 2) and ask it to review:\\
      {\small ``Read the analysis script and identify any bugs, silent data loss, or decisions that could change the results.''}
    \item Agent 2 has fresh context and no attachment to the code
  \end{enumerate}

  \vspace{0.5em}
  \textbf{Why this is powerful:}
  \begin{itemize}
    \item The writing agent has blind spots --- it wrote the code and ``believes'' it is correct
    \item A fresh agent reads the code like a skeptical reviewer
    \item Catches issues you would miss because you watched it being built
  \end{itemize}

  \vspace{0.5em}
  \alert{Treat your agents like co-authors who check each other's work, not a single assistant you trust unconditionally.}
\end{frame}

% ══════════════════════════════════════════════════════════════════
%  STEP 5 AND CONTEXT WINDOW
% ══════════════════════════════════════════════════════════════════

\begin{frame}{Step 4: Save or Re-instruct}
  \textbf{If the output passes verification:} keep it. Claude Code saves its changes to your files after each step, so the results are already on disk.

  \vspace{0.5em}
  \textbf{If something is wrong:} identify the specific error and re-instruct. Do not ask Claude Code to ``fix it'' --- tell it \textit{what} to fix and \textit{how}.

  \vspace{0.5em}
  In Session 3 we will learn \textbf{git} --- a tool that lets you snapshot your project before each change and revert if anything goes wrong. For now, know that the save step exists and that there is a better way to manage it coming soon.

  \vspace{0.5em}
  \begin{block}{The full cycle}
    Describe $\rightarrow$ Review plan $\rightarrow$ Generate \& validate $\rightarrow$ Save or re-instruct\\[0.3em]
    Repeat for each task in your analysis. One cycle per task, not one cycle for the whole project.
  \end{block}
\end{frame}

% ══════════════════════════════════════════════════════════════════
%  WHAT TO DELEGATE
% ══════════════════════════════════════════════════════════════════

\begin{frame}{What to Delegate, What to Verify, What to Keep}
  \begin{columns}[T]
    \begin{column}{0.32\textwidth}
      {\color{green!50!black}\textbf{Delegate freely}}\\[0.3em]
      {\small
      Boilerplate \& docs\\
      File management\textsuperscript{\dag}\\
      Reformatting code}
    \end{column}
    \begin{column}{0.32\textwidth}
      {\color{accent}\textbf{Delegate + verify}}\\[0.3em]
      {\small
      Data loading \& cleaning\\
      Data merges\\
      Variable construction\\
      Summary statistics\\
      Robustness variants\\
      Visualization\\
      \LaTeX{} tables}
    \end{column}
    \begin{column}{0.32\textwidth}
      {\color{red!70!black}\textbf{Do not delegate}}\\[0.3em]
      {\small
      Choosing specifications\\
      Interpreting results\\
      Identification strategy\\
      Sample selection\\
      Causal claims}
    \end{column}
  \end{columns}

  \vspace{0.5em}
  {\scriptsize \textsuperscript{\dag} Be careful with file management --- Claude Code can delete files. Always review before approving.}

  \vspace{0.5em}
  The test: \textit{If a referee questioned this choice, who would need to answer --- you or the RA?}

  \vspace{0.5em}
  Today's exercise focuses on the middle column: tasks where Claude Code is productive but where you must verify every output.
\end{frame}

% ══════════════════════════════════════════════════════════════════
%  EXERCISE SETUP
% ══════════════════════════════════════════════════════════════════

\begin{frame}{Today's Data: The Swiss Franc Shock}
  \textbf{Background:} In January 2015, the Swiss National Bank removed its currency peg.
  The Swiss franc surged against the Polish zloty. Polish households with CHF-denominated
  mortgages saw their payments spike overnight.

  \vspace{0.5em}
  \textbf{The survey:} 2,044 Polish adults, interviewed October 2015. Respondents were
  randomly assigned to receive different information about the crisis before being asked
  whether the government should intervene to help affected borrowers.

  \vspace{0.5em}
  \textbf{Your files:}
  \begin{itemize}
    \item \texttt{swiss\_franc\_survey.csv} --- 2,044 respondents, 11 variables
    \item \texttt{respondent\_demographics.csv} --- income and ideology data (more rows than the survey)
    \item \texttt{codebook.md} --- variable descriptions
  \end{itemize}

  \vspace{0.3em}
  {\scriptsize Source: Ahlquist, Copelovitch, and Walter, ``The Political Consequences of External Economic Shocks: Evidence from Poland,'' \textit{American Journal of Political Science} (2020).}
\end{frame}

\begin{frame}{The Research Question}
  \centering
  \vspace{1em}
  {\Large Does providing information about the Swiss franc crisis}\\[0.3em]
  {\Large increase support for government intervention}\\[0.3em]
  {\Large to help affected borrowers?}

  \vspace{1.5em}
  \begin{columns}[T]
    \begin{column}{0.45\textwidth}
      \textbf{Treatment (randomized):}
      \begin{itemize}
        \item Control (no info)
        \item Factual information
        \item Historical context
        \item Hungary's policy response
      \end{itemize}
    \end{column}
    \begin{column}{0.45\textwidth}
      \textbf{Outcome:}
      \begin{itemize}
        \item Supports government intervention (binary)
      \end{itemize}
      \vspace{0.5em}
      \textbf{Moderator:}
      \begin{itemize}
        \item FX loan exposure status
      \end{itemize}
    \end{column}
  \end{columns}

  \vspace{1em}
  The methods are straightforward. The skill is staying in control of the workflow.
\end{frame}

% ══════════════════════════════════════════════════════════════════
%  EXERCISE PART 1: MERGE AND EXPLORE
% ══════════════════════════════════════════════════════════════════

\begin{frame}{Exercise Part 1: Merge and Explore (30 min)}
  \textbf{Use plan mode for every step. Review every plan before approving.}

  \vspace{0.8em}
  \begin{enumerate}
    \item \textbf{Merge} the demographics data into the survey.
    \begin{itemize}
      \item The demographics file has more rows than the survey
      \item \alert{You decide the join type.} Do not let Claude Code choose.
    \end{itemize}

    \vspace{0.3em}
    \item \textbf{Explore} the merged data: summary statistics, treatment group sizes, outcome distribution.
    \begin{itemize}
      \item \alert{Before you look at results: write down what you expect to see.}
    \end{itemize}

    \vspace{0.3em}
    \item \textbf{Verify} the merge.
    \begin{itemize}
      \item What should you check? Write down your verification checks \textit{before} running them.
      \item Then run them.
    \end{itemize}
  \end{enumerate}
\end{frame}

\begin{frame}{Part 1 Debrief: What Did You Verify?}
  \textbf{What did you check?}

  \vspace{1.5em}
  \pause
  \textbf{You should have verified:}
  \begin{itemize}
    \item Row count: exactly 2,044 rows (no rows gained or lost)
    \item Join type: only a left join gives the correct result
    \item No duplicate respondent IDs in the output
    \item Treatment groups: 480 / 549 / 520 / 495
    \item Proportion supporting intervention: $\sim$0.479 (unweighted)
    \item FX exposure: 1,918 / 69 / 55
  \end{itemize}

  \vspace{0.5em}
  Did your numbers match? Did Claude Code choose the right join type?
\end{frame}

% ══════════════════════════════════════════════════════════════════
%  EXERCISE PART 2: SPECIFY AND ESTIMATE
% ══════════════════════════════════════════════════════════════════

\begin{frame}{Exercise Part 2: Specify and Estimate (30 min)}
  \begin{enumerate}
    \setcounter{enumi}{3}
    \item \textbf{Base model:} Logistic regression of \texttt{supports\_intervention} on \texttt{treatment}.
    \begin{itemize}
      \item Use survey weights
      \item What should the reference category be? Tell Claude Code explicitly.
    \end{itemize}

    \vspace{0.3em}
    \item \textbf{Add FX exposure.} Did the sample size change? Why or why not?

    \vspace{0.3em}
    \item \textbf{Add demographic controls:} age, female, education, urban/rural, income, left-right.
    \begin{itemize}
      \item \alert{Before running: predict how many observations you will have.}
      \item Then check. Were you right?
    \end{itemize}

    \vspace{0.3em}
    \item \textbf{Add an interaction} between treatment and FX exposure.
  \end{enumerate}
\end{frame}

\begin{frame}{Part 2 Debrief: What Happened to Your Sample?}
  \textbf{How many observations did you have in each model?}

  \vspace{1em}
  \pause
  \begin{alertblock}{The Missing Data Trap}
    Adding control variables can silently change your sample.
  \end{alertblock}

  \vspace{0.3em}
  \begin{tabular}{lrl}
    \toprule
    \textbf{Specification} & \textbf{N} & \textbf{What happened} \\
    \midrule
    Treatment only & 2,036 & Full sample (8 missing outcome) \\
    + FX exposure & 2,035 & 2 missing FX status \\
    + age, female, education, urban/rural & 2,035 & No additional missingness \\
    + income quintile & $\sim$1,473 & \color{red!70!black}{565 missing income} \\
    + left-right scale & $\sim$1,499 & \color{red!70!black}{542 missing left-right} \\
    + both & $\sim$1,151 & \color{red!70!black}{44\% of sample dropped} \\
    \bottomrule
  \end{tabular}

  \vspace{0.5em}
  \textbf{Did Claude Code mention this when it added the controls?}
\end{frame}

% ══════════════════════════════════════════════════════════════════
%  EXERCISE PART 3: ROBUSTNESS AND LATEX
% ══════════════════════════════════════════════════════════════════

\begin{frame}{Exercise Part 3: Robustness and Output (30 min)}
  \begin{enumerate}
    \setcounter{enumi}{7}
    \item \textbf{Generate robustness checks} from your base model:
    \begin{itemize}
      \item With and without survey weights
      \item Different control sets (track how N changes each time!)
      \item Collapse three treatment arms into ``any information'' vs.\ control
      \item Restrict to complete cases so N is constant across specs
    \end{itemize}

    \vspace{0.5em}
    \item \textbf{Format as \LaTeX.} Ask Claude Code to produce a publication-style
    regression table with all specifications side by side. Compile to PDF.

    \vspace{0.5em}
    \item \textbf{Commit.} Save all scripts, output, and the \LaTeX{} table.
  \end{enumerate}
\end{frame}

\begin{frame}{The Robustness Workflow}
  \centering
  {\large You specify the core analysis.\\
  Claude Code generates the variations.\\
  You verify each one.}

  \vspace{1.5em}
  \textbf{Claude Code never chooses what robustness checks to run.}

  \vspace{1em}
  \textbf{For each specification, check:}
  \begin{itemize}
    \item Does it use the right sample? (What is N?)
    \item Does the code match what you asked for?
    \item Are the coefficients in a plausible range?
    \item Do the signs on control variables make sense?
  \end{itemize}

  \vspace{0.5em}
  \begin{alertblock}{The common mistake}
    A robustness check that uses a different sample than the main specification
    is not a robustness check --- it is a different analysis.
  \end{alertblock}
\end{frame}

\begin{frame}{\LaTeX{} Tables: Bringing Results Into Your Paper}
  Until now, your results are printed to the console or saved as CSVs.\\
  In your actual workflow, they need to end up in a paper.

  \vspace{0.8em}
  \textbf{Why this matters:}
  \begin{itemize}
    \item Claude Code can write \LaTeX{} tables directly from regression output
    \item The table lives in your project directory, under version control
    \item When you update the analysis, Claude Code updates the table
    \item No manual reformatting, no copy-paste errors
  \end{itemize}

  \vspace{0.8em}
  \textbf{Today:} generate one table. In Session 3, we bring the full paper into VS Code.

  \vspace{0.5em}
  \begin{block}{The bigger idea}
    The more of your research that lives inside the project directory, the more Claude Code
    can operate on. Code, data, output, tables --- and eventually the paper itself.
  \end{block}
\end{frame}

% ══════════════════════════════════════════════════════════════════
%  WHAT TO WATCH FOR
% ══════════════════════════════════════════════════════════════════

\begin{frame}{What to Watch For During the Exercise}
  This exercise is designed to surface moments where Claude Code makes choices for you.

  \vspace{0.8em}
  \begin{itemize}
    \item \textbf{The merge:} Did Claude Code pick the right join type? Did it mention the extra rows?
    \item \textbf{Reference categories:} Did it set treatment to ``control'' or pick alphabetically?
    \item \textbf{Missing data:} When you added controls, how many observations did you lose? Did Claude Code mention this?
    \item \textbf{Survey weights:} Did it use them unless you told it to?
    \item \textbf{Income = 0:} ``No income'' is coded as 0. How did Claude Code treat it?
  \end{itemize}

  \vspace{0.8em}
  \alert{These are the moments where your judgment matters.}\\
  The workflow is: notice the decision, direct Claude Code explicitly, then verify the output.
\end{frame}

% ══════════════════════════════════════════════════════════════════
%  DEBRIEF
% ══════════════════════════════════════════════════════════════════

\begin{frame}{Debrief}
  \textbf{Discussion questions:}

  \vspace{0.8em}
  \begin{enumerate}
    \item Did your merge lose any observations? How did you handle weekends?
    \item How many observations did you lose when you added controls?\\Did Claude Code tell you?
    \item How many of your robustness specifications were correct on the first try?
    \item What did you catch in the plan that you would have missed in the output?
  \end{enumerate}

  \vspace{1em}
  \begin{block}{Key takeaway}
    The assertions and the plan review are not extra work --- they are the workflow. The 30 seconds it takes to read a plan or check a sample size can save hours of debugging downstream.
  \end{block}
\end{frame}

\begin{frame}{Next Session}
  \textbf{Session 3: Code Auditing, Git, and \LaTeX{} in VS Code}

  \vspace{1em}
  \begin{itemize}
    \item Claude Code reads and critiques \textit{your} code instead of writing it
    \item Git as your safety net: commit $\rightarrow$ instruct $\rightarrow$ diff $\rightarrow$ revert
    \item Bringing your paper into the project: \LaTeX{} in VS Code
  \end{itemize}

  \vspace{1em}
  \textbf{Before next session:}
  \begin{itemize}
    \item Try the core workflow on something from your own research
    \item Note where Claude Code made a decision you did not ask for
    \item Install the \LaTeX{} Workshop extension in VS Code (instructions in setup guide)
  \end{itemize}
\end{frame}

\begin{frame}
\end{frame}

\end{document}
